\documentclass[12pt]{article}

% Use algorithm and algpseudocode packages so \State and related commands are available

\newcommand{\duedate}{---}
\newcommand{\assignment}{Provable-Guarantees (Lec6) Summary} % Change to "Problem Set X"

% Change the following to your name and UNI.
\newcommand{\name}{Aksel Kretsinger-Walters, adk2164}
\newcommand{\email}{adk2164@columbia.edu}

% NOTE: Defining collaborators is optional; to not list collaborators, comment out the
% line below. Maximum of two collaborators per problem set.
%\newcommand{\collaborators}{Vig Vigerton (\texttt{UNI}), Alice Bob (\texttt{UNI})}
% No collaborators on PS 0

\makeatletter
\def\input@path{{../}{../../}{../../../}} % add as many parents as you need
\makeatother
\input{pset_template_RL.tex} %% DO NOT CHANGE THIS LINE

\begin{document}
\psetheader %% DO NOT CHANGE THIS LINE

\section*{Lecture 6 Summary: Provable Guarantees for Policy Gradient Methods}

\subsection*{Overview}
This lecture studies the theoretical foundations and convergence guarantees of policy
gradient methods, based on \textit{Kakade \& Langford (2002)} and \textit{Schulman et al.
(2015)}. It highlights situations where naive policy gradient updates fail, and presents
provably efficient variants such as the \emph{Conservative Greedy Algorithm} and
\emph{Trust Region Policy Optimization (TRPO)}.

\subsection*{1. When Policy Gradients Fail}

\paragraph{Example 1 – Exponential Trajectories.}
In a chain MDP where each state $s \in \{n,\dots,1\}$ has two forward actions and one
backward action, the uniform policy’s probability of reaching the rewarding state 1
decays exponentially with $n$.
Hence, the expected gradient magnitude
\[
		\|\nabla_\theta \rho(\pi)\| \le \frac{n}{1-\gamma}\left(\frac{1}{2}\right)^{n/2},
\]
implying exponentially many samples are needed to obtain a nonzero gradient estimate.

\paragraph{Example 2 – Two-State Trap.}
Even with two states $(i,j)$, policy gradients may worsen the policy.
Because gradient updates scale with the stationary distribution $d^\pi(s)$,
the high-probability state $i$ dominates learning.
The update
\[
		\theta_{s,a} \leftarrow \theta_{s,a} + \alpha\,
		d^\pi(s)\pi_\theta(s,a)\big(Q^\pi(s,a)-V^\pi(s)\big)
\]
can increase the probability of looping on $i$, worsening performance before eventual
recovery in exponential time.

\subsection*{2. Performance Difference Lemma}

A key result that underlies all policy improvement analysis:
\[
		\rho(\tilde\pi)-\rho(\pi)
		= \frac{1}{1-\gamma}\,
		\mathbb{E}_{s\sim(1-\gamma)d^{\tilde\pi},\,a\sim\tilde\pi}
		[A^\pi(s,a)],
		\quad A^\pi(s,a)=Q^\pi(s,a)-V^\pi(s).
\]
This identity expresses the performance gap between any two policies in terms of the
\emph{advantage} of $\pi$ under the new policy’s state distribution.

\subsection*{3. Conservative Greedy Policy Improvement}

To ensure monotonic improvement, \textit{Kakade \& Langford (2002)} propose the update:
\[
		\pi_{\text{new}}(s,a) = (1-\alpha)\pi(s,a) + \alpha\pi'(s,a),
\]
where $\pi'$ is a candidate (greedy) policy and $\alpha$ controls conservativeness.

\paragraph{Guaranteed Improvement.}
For any $\pi'$,
\[
		\rho(\pi_{\text{new}}) - \rho(\pi)
		\ge \alpha\,A_\pi(\pi') - \frac{2\alpha^2\gamma\epsilon}{1-\gamma(1-\alpha)},
\]
where
\[
		A_\pi(\pi') = \tfrac{1}{1-\gamma}\mathbb{E}_{s,a\sim d^\pi,\pi'}[A_\pi(s,a)],\quad
		\epsilon = \tfrac{1}{1-\gamma}\max_s\big|\!\sum_a \pi'(s,a)A_\pi(s,a)\big|.
\]

\paragraph{Choosing Parameters.}
If rewards are bounded by $R$, selecting
\[
		\alpha = \frac{A_\pi(\pi')(1-\gamma)^3}{4R}
\]
yields the provable bound
\[
		\rho(\pi_{\text{new}}) - \rho(\pi)
		\ge \frac{(1-\gamma)^3A_\pi(\pi')^2}{8R}.
\]
Thus each update is guaranteed non-negative improvement.

\paragraph{Termination Guarantee.}
The conservative greedy algorithm stops after at most
\[
		O\!\left(\frac{R^2}{\delta^2}\right)
\]
iterations to find a $\pi$ satisfying
\[
		(1-\gamma)\max_{\pi'}A_\pi(\pi')\le \delta.
\]

\subsection*{4. Local vs. Global Optimality}

Even a locally optimal policy can differ significantly from the global optimum if state
distributions differ.
\[
		\rho(\pi^*;\mu^*) - \rho(\pi;\mu^*)
		\le \frac{\delta}{(1-\gamma)^2}
		\left\|\frac{d_{\pi^*,\mu^*}}{d_{\pi,\mu}}\right\|_\infty.
\]
Choosing a more uniform initial state distribution mitigates this gap.

\subsection*{5. Trust Region Policy Optimization (TRPO)}

\textit{Schulman et al. (2015)} generalized the conservative update by replacing the
mixture constraint with a \emph{trust region} constraint based on total-variation or KL divergence:
\[
		\rho(\tilde\pi)-\rho(\pi)
		\ge A_\pi(\tilde\pi)
		- \frac{2\gamma\epsilon}{1-\gamma}\, D^{\max}_{TV}(\pi,\tilde\pi)^2
		\ge A_\pi(\tilde\pi)
		- \frac{4\gamma\epsilon}{1-\gamma}\, D^{\max}_{KL}(\pi,\tilde\pi).
\]
This ensures safe updates without forming explicit policy mixtures.

\paragraph{TRPO Objective.}
Find the next policy parameter $\tilde\theta$ by solving
\[
		\max_{\tilde\theta}\; A_\theta(\tilde\theta)
		\quad\text{s.t.}\quad
		D^{\max}_{KL}\big(\pi_\theta\,\|\,\pi_{\tilde\theta}\big)\le \eta,
\]
where $\eta$ controls the trust-region size.

\subsection*{6. Key Takeaways}
\begin{itemize}
  \item Naive policy gradients can converge exponentially slowly or even worsen performance.
		\item The \textbf{Performance Difference Lemma} formalizes exact performance change
				between policies.
  \item \textbf{Conservative Greedy Algorithm} ensures monotonic improvement via step-size control.
  \item \textbf{TRPO} removes mixtures by constraining updates through KL or TV divergence.
		\item Both frameworks form the theoretical foundation for modern stable policy-gradient
				methods (e.g., PPO).
\end{itemize}

\end{document}
